% !TEX root = ../main.tex

\begin{acknowledgments}
  % 学位论文正文和附录之后，一般应放置致谢（或后记），
  % 主要感谢导师和对论文工作有直接贡献和帮助的人士和单位。
  % 包括：
  % a) 国家科学基金，资助研究工作的奖学金基金，合同单位，资助或支持的企业、组织或个人；
  % b) 协助完成研究工作和提供便利条件的组织或个人；
  % c) 在研究工作中提出建议和提供帮助的人；
  % d) 给予转载和引用权的资料、图片、文献、研究思想和设想的所有者；
  % e) 其他应感谢的组织和个人。
  % 致谢言语应谦虚诚恳，实事求是。
  % 字数不超过1000个汉字，格式要求与正文同，标题“致谢”的格式要求与章标题同。

  不知不觉在科大读研已七年，七年时间转瞬即逝，在这七年时间里也学会了很多，
  在此我要对指导过我的老师，帮助过我的同学，照顾过我的朋友说一声“谢谢”！

  首先感谢我的导师***，为我提供了读研的宝贵机会，
  使我能够继续深造并且很幸运地在科研的道路上前行。
  此论文是在导师的指导下完成的，他们在实验结果的讨论、论文的结构、
  撰写和修改等方面都提出了宝贵的意见，也提供了很多的帮助。
  导师深厚的学术功底，敏锐的科学洞察力，独到的见解使我受益良多，
  也会深刻影响我以后的工作和生活。
  在此，向他们致以最诚挚的感谢！

\end{acknowledgments}
